% =====================================================================
%  Instrucciones para evaluadores independientes
%  Proyecto SIMPA -- Equipo AHMRV -- ISR-401 -- UTEQ -- 2026-2027 PPA
%  Instrumento del cuasi-experimento del Enfoque 1 -- EXP-06
%  Compilacion: pdflatex instrucciones_evaluador.tex  (x2)
% =====================================================================
\documentclass[11pt,a4paper]{article}

\usepackage[utf8]{inputenc}
\usepackage[T1]{fontenc}
\usepackage[margin=2.3cm]{geometry}
\usepackage{array}
\usepackage{booktabs}
\usepackage{longtable}
\usepackage[table]{xcolor}
\usepackage{enumitem}
\usepackage{tcolorbox}
\usepackage{fancyhdr}
\usepackage{titlesec}
\usepackage{lastpage}
\usepackage[hidelinks]{hyperref}

\definecolor{uteqverde}{HTML}{1F6F3C}
\definecolor{grisclaro}{HTML}{F2F2F2}

\newcolumntype{L}[1]{>{\raggedright\arraybackslash}p{#1}}
\newcolumntype{C}[1]{>{\centering\arraybackslash}p{#1}}

\titleformat{\section}{\normalfont\large\bfseries\color{uteqverde}}{\thesection.}{0.5em}{}
\titlespacing*{\section}{0pt}{12pt}{5pt}
\titleformat{\subsection}{\normalfont\normalsize\bfseries}{\thesubsection}{0.5em}{}

\pagestyle{fancy}
\fancyhf{}
\renewcommand{\headrulewidth}{0.4pt}
\lhead{\footnotesize\color{uteqverde}Instrucciones para evaluadores independientes}
\rhead{\footnotesize SIMPA --- Equipo AHMRV}
\cfoot{\footnotesize Página \thepage\ de \pageref{LastPage}}

\setlength{\parskip}{4pt}
\setlength{\parindent}{0pt}

\begin{document}

\begin{center}
{\footnotesize UNIVERSIDAD TÉCNICA ESTATAL DE QUEVEDO\\
Facultad de Ciencias de la Computación --- Carrera de Software}\\[10pt]
{\Large\bfseries\color{uteqverde} INSTRUCCIONES PARA EVALUADORES INDEPENDIENTES}\\[3pt]
{\large Cuasi-experimento de calidad de requisitos funcionales}\\[2pt]
{\footnotesize Instrumento del estudio registrado en OSF \texttt{4z35d}}
\end{center}

\vspace{6pt}

\section{Qué se le pide}

Va a recibir un archivo (\texttt{requisitos\_cegados.csv}) con 50 requisitos
funcionales, identificados solo con un código \texttt{R-001} a
\texttt{R-050}, en orden aleatorio. Para cada uno, debe asignar una
puntuación de 1 a 5 en cada una de las cinco dimensiones de la rúbrica
adjunta (\texttt{rubrica\_evaluacion.pdf}), registrando el resultado en la
hoja de puntuación (\texttt{hoja\_puntuacion.pdf}).

\section{Lo que necesita saber antes de empezar}

\begin{tcolorbox}[colback=grisclaro, colframe=uteqverde, boxrule=0.6pt]
\textbf{Los 50 requisitos provienen de dos orígenes distintos: la mitad
fueron elaborados por analistas humanos y la otra mitad por un modelo de
lenguaje.} El archivo que usted recibe \textbf{no indica cuál es cuál}, y
esa información no debe buscarse ni preguntarse durante la evaluación.
Evaluar sin saber el origen es la condición que hace válida la
comparación.
\end{tcolorbox}

\begin{itemize}[leftmargin=1.1cm]
  \item No intente adivinar ni comentar con el equipo cuál cree que es el
        origen de un requisito mientras evalúa.
  \item Evalúe cada requisito únicamente por su propio contenido: claridad,
        completitud, verificabilidad y consistencia interna, conforme a la
        rúbrica.
  \item Puntúe los 50 requisitos. No omita ninguno ni los evalúe en un
        orden distinto al que aparecen en el archivo.
  \item Trabaje de forma individual. No consulte su puntuación con los
        otros dos evaluadores antes de entregar su hoja de puntuación
        completa.
  \item Si algún requisito le resulta ambiguo o incompleto, puntúelo según
        la rúbrica igualmente y anote la ambigüedad en el campo de
        observaciones de la hoja de puntuación --- no lo descarte ni lo
        deje en blanco.
\end{itemize}

\section{Materiales que debe tener}

\begin{longtable}{|L{5.5cm}|L{9.0cm}|}
\hline
\rowcolor{grisclaro}\textbf{Archivo} & \textbf{Para qué} \\ \hline
\endhead
\texttt{requisitos\_cegados.csv} & Los 50 requisitos a evaluar, sin pista de origen \\ \hline
\texttt{rubrica\_evaluacion.pdf} & Los criterios y la escala de 1 a 5 para cada dimensión \\ \hline
\texttt{hoja\_puntuacion.pdf} & Donde registra su puntuación de cada requisito \\ \hline
\texttt{consentimiento\_evaluador.pdf} & Debe estar firmado antes de recibir los demás archivos \\ \hline
\end{longtable}

\section{Procedimiento}

\begin{enumerate}[label=\arabic*.]
  \item Firma del consentimiento informado.
  \item Recepción de esta instrucción, la rúbrica y el archivo cegado.
  \item Lectura de la rúbrica y resolución de dudas sobre los criterios
        \textbf{antes} de ver los requisitos (las dudas sobre un requisito
        específico ya no se resuelven una vez iniciada la evaluación, para
        no introducir sesgo).
  \item Evaluación individual e independiente de los 50 requisitos.
  \item Entrega de la hoja de puntuación completa al control documental
        del experimento.
\end{enumerate}

\section{Confidencialidad}

Su identidad como evaluador o evaluadora se registra únicamente con un
código (\texttt{EV-1}, \texttt{EV-2} o \texttt{EV-3}) en los archivos
públicos del proyecto. El detalle de su identidad y el consentimiento
firmado se conservan de forma restringida, conforme al mismo protocolo de
privacidad aplicado al resto de la evidencia del proyecto.

\vfill
\hrule
\vspace{0.2cm}
{\small Cualquier duda sobre este procedimiento, antes de comenzar, debe
dirigirse al control documental del experimento (Denisses Huilcapi), no a
la persona que le entregó los materiales.\par}

\end{document}
